\PuzzleChapterIntro{The Missing Treasure}{Warmup \textbullet\ Base Rates \textbullet\ Posterior}{This is a two-stage forensic puzzle: first, you must model a state-dependent sampling process (the ``sock ritual'') to derive a precise base rate, and then you must apply that prior to a Bayesian likelihood comparison to identify which suspect possesses guilty knowledge.}
\Lore{

They called it treasure only because it felt like one: velvet-wrapped, rain-chilled, still gritty at the seams.

Every drawer received a brittle one-use strip, but the drawer that mattered received a strip chosen by chance---because chance is easier to argue with than memory.

}

\Problem

Five friends---\textbf{Raven}, \textbf{Athena}, \textbf{Quill}, \textbf{Kestrel}, and \textbf{Sable}---found a velvet-wrapped treasure while out walking.

They agreed Raven would keep it overnight for safekeeping. Raven left the group at \textbf{6:00 PM} and put the treasure in one dresser drawer.

\medskip

\noindent Raven sealed \emph{every} dresser drawer with a one-use security strip.
(A strip irreversibly shatters if that drawer is opened even slightly.)

\medskip

\noindent\textbf{The hidden-mark strip.}
Only the treasure drawer got a special strip whose underside fold carried \emph{one} hidden mark:
either a small \textbf{Wave} or a small \textbf{Knot}. The mark is invisible unless the strip shatters.

Which mark Raven used was decided by a ritual performed \emph{before} sealing:

\begin{quote}
Raven has a bag containing \textbf{5 pairs of socks}, each pair a different color (so 10 socks total).
She draws \textbf{4 single socks} at random and holds them aside.
If the 4 socks in hand contain any \emph{complete color-pair} (both socks of some color),
Raven discards that complete pair and immediately draws replacements (one sock per discarded sock) from the bag
until 4 socks are in hand again.

She repeats this discard-and-replace process until either:
\begin{itemize}
  \item the hand contains \textbf{4 socks of 4 different colors} (then a \textbf{Wave} strip is used), or
  \item the bag becomes \textbf{empty} before 4 different colors are ever reached (then a \textbf{Knot} strip is used).
\end{itemize}
\end{quote}

\medskip

At \textbf{8:10 PM}, Raven checked the dresser in passing: all strips were intact.

At \textbf{8:40 PM}, Raven checked again and found \emph{exactly one} strip shattered.
The treasure was gone. The shattered strip's hidden mark was a \textbf{Wave}.

\medskip

Raven did \emph{not} tell anyone which mark appeared. The next day Raven questioned the other four friends.
To keep the stories from changing midstream, Raven did the following:

\begin{itemize}
\item She asked each person, alone, to write on a slip of paper:
\textit{``Wave or Knot: which do you believe was under the shattered strip?''}
\item Each person sealed their slip in wax \emph{before} speaking further.
\item Everyone's slip and statement below is \textbf{truthful}:
if someone actually opened the drawer and saw the mark, they wrote what they saw;
otherwise they wrote according to their usual habit described below.
\item Raven had \textbf{no reason beforehand} to suspect one friend more than another.
\end{itemize}

When Raven opened the sealed slips, the slips read:

\begin{quote}
\textbf{Athena's slip:} \textbf{Wave}. \\
\textbf{Athena added:} ``If I don't know, I always bet on whichever mark is \emph{more likely} under your sock ritual.''
\end{quote}

\begin{quote}
\textbf{Kestrel's slip:} \textbf{Knot}. \\
\textbf{Kestrel added:} ``If I don't know, I always bet on whichever mark is \emph{less likely}. I like long shots.''
\end{quote}

\begin{quote}
\textbf{Sable's slip:} \textbf{Wave}. \\
\textbf{Sable added:} ``If I don't know, I try to `match the odds': I choose Wave with the \emph{same frequency}
that your ritual produces Wave.''
\end{quote}

\begin{quote}
\textbf{Quill's slip:} \textbf{Wave}. \\
\textbf{Quill added:} ``If I don't know, I flip a fair coin: Wave on heads, Knot on tails.''
\end{quote}

\VaultDirective{
Identify the most likely thief.
}


\SolutionPage

\textbf{1. The Evidence}
The shattered strip revealed a \textbf{Wave}.
\begin{itemize}
    \item \textbf{The Thief:} Must write \textbf{Wave} (because they saw it and are compelled to be truthful about what they saw).
    \item \textbf{The Innocent:} Writes Wave or Knot based on their guessing strategy.
\end{itemize}
\textbf{Kestrel} wrote Knot. Since the thief would have written Wave, Kestrel is \textbf{innocent}.
We must determine who among Athena, Sable, and Quill is most likely to be the thief.

\medskip

\textbf{2. The Base Rate (Sock Ritual)}
We calculate \(p = P(\text{Wave})\), the probability of drawing 4 distinct colors.
Total socks: \(10\) (\(5\) pairs). Initial draw: \(4\) socks. Total combinations: \(\binom{10}{4} = 210\).

\begin{itemize}
    \item \textbf{Immediate Wave (4 different colors):}
    Choose 4 pairs from 5, then 1 sock from each.
    \[ P(\text{Wave}_1) = \frac{\binom{5}{4} \times 2^4}{210} = \frac{5 \times 16}{210} = \frac{8}{21} \]
    \item \textbf{One Pair (2 different colors + 1 pair):}
    Choose 1 pair to match, 2 pairs for the singles.
    \[ P(\text{1 Pair}) = \frac{\binom{5}{1} \times \binom{4}{2} \times 2^2}{210} = \frac{5 \times 6 \times 4}{210} = \frac{4}{7} \]
    \textit{Action:} Discard the pair. We hold 2 different socks. Bag has 6 socks left (contains the mates for our held socks + 2 fresh pairs).
    To get a Wave, we must draw the two fresh colors (one from each new pair) from the 6 remaining socks.
    \[ P(\text{Success} \mid \text{1 Pair}) = \frac{2 \times 2}{\binom{6}{2}} = \frac{4}{15} \]
\end{itemize}
\[
p = \frac{8}{21} + \left( \frac{4}{7} \times \frac{4}{15} \right) = \frac{40}{105} + \frac{16}{105} = \frac{56}{105} = \frac{8}{15}.
\]
Since \(8/15 > 7/15\), Wave is the \textbf{more likely} outcome.

\medskip

\textbf{3. Suspect Probabilities}
We compare the likelihood of the observed slips under each suspect's guilt.
\begin{itemize}
    \item \textbf{Athena (Guilty):} She writes Wave (thief). Sable guesses Wave (\(p=8/15\)). Quill guesses Wave (\(p=1/2\)).
    \[ L_A = 1 \times \frac{8}{15} \times \frac{1}{2} = \frac{4}{15} \approx 0.26 \]
    \item \textbf{Sable (Guilty):} She writes Wave (thief). Athena guesses Wave (bets ``more likely'', \(p=1\)). Quill guesses Wave (\(p=1/2\)).
    \[ L_S = 1 \times 1 \times \frac{1}{2} = \frac{1}{2} = 0.50 \]
    \item \textbf{Quill (Guilty):} He writes Wave (thief). Athena guesses Wave (\(p=1\)). Sable guesses Wave (\(p=8/15\)).
    \[ L_Q = 1 \times 1 \times \frac{8}{15} = \frac{8}{15} \approx 0.53 \]
\end{itemize}
The likelihood is highest for Quill.

\FinalAnswer{QUILL}

\NextPage